\begin{example}[Projective zero-space]\label{ex:vi36-01}
\(\mathbb P^0_k=\Proj k[x_0]\) has one point and is isomorphic to \(\Spec k\).
\end{example}

\begin{example}[A point in two charts]\label{ex:vi36-02}
The point \([1:2:3]\in\mathbb P^2_k\) lies in all three standard charts when \(2,3\ne0\).  In \(U_0\) its coordinates are \((2,3)\); in \(U_1\) they are \((1/2,3/2)\).
\end{example}

\begin{example}[A point at infinity]\label{ex:vi36-03}
Relative to the affine chart \(U_0\subset\mathbb P^2\), the point \([0:1:2]\) lies on the line at infinity \(x_0=0\).
\end{example}

\begin{example}[Projective line transition]\label{ex:vi36-04}
On \(\mathbb P^1\), \(t=x_1/x_0\) and \(u=x_0/x_1\) satisfy \(u=t^{-1}\) on \(\mathbb G_m\).
\end{example}

\begin{example}[Dehomogenizing a line]\label{ex:vi36-05}
The homogeneous linear equation \(x_0+x_1+x_2=0\) becomes \(1+t_1+t_2=0\) on \(U_0\).
\end{example}

\begin{example}[Dehomogenizing a conic]\label{ex:vi36-06}
The equation \(x_0x_2-x_1^2=0\) becomes \(t_2-t_1^2=0\) on \(U_0\), displaying the affine parabola chart of the projective conic.
\end{example}

\begin{example}[The line at infinity of a plane curve]\label{ex:vi36-07}
If an affine polynomial \(f(x,y)\) is homogenized to \(F(x_0,x_1,x_2)\), its points at infinity are detected set-theoretically by \(F(0,x_1,x_2)=0\).
\end{example}

\begin{example}[A projective coordinate change]\label{ex:vi36-08}
An invertible matrix \(A\in\mathrm{GL}_{n+1}(k)\) sends a nonzero vector to a nonzero vector and therefore induces a bijection \([v]\mapsto[Av]\) on \(\mathbb P^n(k)\).  The corresponding scheme automorphism will be interpreted systematically in later projective-morphism language.
\end{example}

\begin{example}[A non-rational real point]\label{ex:vi36-09}
The closed point of \(\mathbb P^1_{\mathbb R}\) defined on \(U_0\) by \(t^2+1\) has residue field \(\mathbb C\).
\end{example}

\begin{example}[Base extension to the complex numbers]\label{ex:vi36-10}
After base changing the preceding point from \(\mathbb R\) to \(\mathbb C\), the polynomial \(t^2+1\) factors and the geometric fiber separates into the two \(\mathbb C\)-points \([1:i]\) and \([1:-i]\).
\end{example}

\begin{example}[A hyperplane]\label{ex:vi36-11}
The locus \(x_n=0\) in \(\mathbb P^n\) has coordinates \([x_0:\cdots:x_{n-1}]\) and is naturally a \(\mathbb P^{n-1}\).
\end{example}

\begin{example}[A projective line in \(\mathbb P^3\)]\label{ex:vi36-12}
The simultaneous linear equations \(x_0=x_1=0\) leave coordinates \([0:0:x_2:x_3]\), giving a copy of \(\mathbb P^1\).
\end{example}

\begin{example}[Sections of \(\mathcal O(1)\)]\label{ex:vi36-13}
\(\Gamma(\mathbb P^n,\mathcal O(1))\) has basis \(x_0,\ldots,x_n\) and dimension \(n+1\).
\end{example}

\begin{example}[Sections of \(\mathcal O(2)\) on \(\mathbb P^2\)]\label{ex:vi36-14}
The six quadratic monomials \(x_0^2,x_0x_1,x_0x_2,x_1^2,x_1x_2,x_2^2\) form a basis of \(\Gamma(\mathbb P^2,\mathcal O(2))\).
\end{example}

\begin{example}[Quadratic Veronese of \(\mathbb P^1\)]\label{ex:vi36-15}
The map \([s:t]\mapsto[s^2:st:t^2]\) lands on the conic \(XZ-Y^2=0\) in \(\mathbb P^2\).
\end{example}

\begin{example}[Segre coordinates for \(\mathbb P^1\times\mathbb P^1\)]\label{ex:vi36-16}
The Segre map sends \(([s:t],[u:v])\) to \([su:sv:tu:tv]\in\mathbb P^3\), whose coordinates satisfy \(z_{00}z_{11}-z_{01}z_{10}=0\).
\end{example}
